\subsection{幺代数中元素的谱}

定义 1. --- 设 A 是 K 上的一个幺代数，e 为其单位元。对每个 \(x \in \mathrm{A}\)，我们把 \(x\) 相对于 A 的谱称为由 \(\lambda \in \mathrm{K}\) 中使 \(\lambda e - x\) 不可逆的那些元素所成的集合。

x 的谱将记作 \(\mathrm{Sp}_{\mathrm{A}}(x)\)，若不致混淆，也记作 \(\mathrm{Sp}(x)\)。\(\mathrm{Sp}_{\mathrm{A}}(x)\) 在 K 中的补集称为 x 的预解集。

注. --- 1) 若 A = \{0\}，则有 Sp(0) = ∅。

\begin{enumerate}
\def\labelenumi{\arabic{enumi})}
\setcounter{enumi}{1}
\item
  若 \(A \neq \{0\}\)，则我们有 \(\mathrm{Sp}(\lambda e) = \{\lambda\}\)，其中 \(\lambda \in K\) 任意。
\item
  要使 \(x \in \mathrm{A}\) 可逆，必须且只需 \(0 \notin \operatorname{Sp}(x)\)。
\item
  设 \(R \in K(X)\) 是一个有理分式，并设 \(x \in A\) 是一个可在 R 中代入的元素，也就是说 (A, IV, p.~20)，存在 P 与 \(Q \in K[X]\)，使得 \(R = P/Q\) 且 \(Q(x)\) 可逆；于是可以构成 A 的元素 \(R(x) = P(x) \cdot Q(x)^{-1} = Q(x)^{-1} \cdot P(x)\)；它不依赖于 P 与 Q 的选取。我们有 \(0 \notin Q(\mathrm{Sp}(x))\)，从而 \(\mathrm{Sp}(x)\) 的每个元素都可在 R 中代入。
\end{enumerate}

我们有 \(\mathrm{R}(\mathrm{Sp}(x)) \subset \mathrm{Sp}(\mathrm{R}(x))\)。事实上，设 \(\lambda \in \mathrm{Sp}(x)\)；存在多项式 \(P_{1}\)，使得 \(\mathrm{R}(\lambda) - \mathrm{R}(\mathrm{X}) = (\lambda - \mathrm{X})\mathrm{P}_{1}(\mathrm{X}) / \mathrm{Q}(\mathrm{X})\)；于是 \(\mathrm{R}(\lambda)e - \mathrm{R}(x) = (\lambda e - x)(\mathrm{P}_{1}(x)/\mathrm{Q}(x))\)，因而 \(\mathrm{R}(\lambda) - \mathrm{R}(x)\) 不可逆，故 \(\mathrm{R}(\lambda) \in \mathrm{Sp}(\mathrm{R}(x))\)。

反之，设域 K 是代数闭的。先设 R 不是常数，我们来证明 \(\mathrm{Sp}(\mathrm{R}(x)) = \mathrm{R}(\mathrm{Sp}(x))\)。设 \(\mu \in \mathrm{Sp}(\mathrm{R}(x))\)。由于 R 不是常数，P - \(\mu Q\) 不是零多项式；设 \(\mu Q - P = \alpha \prod (\lambda_i - X)\) 是一个一次因式的分解，于是 \(\mu e - R(x) = \alpha \prod (\lambda_i e - x) Q(x)^{-1}\)。由于 \(\mu e - R(x)\) 不可逆，存在 i 使 \(\lambda_i e - x\) 不可逆，因而 \(\lambda_i \in \mathrm{Sp}(x)\)，于是 \(\mathrm{R}(\lambda_i) = \mu \in \mathrm{R}(\mathrm{Sp}(x))\)。

当 R 为常数时，等式 \(\mathrm{Sp}(\mathrm{R}(x)) = \mathrm{R}(\mathrm{Sp}(x))\) 同样成立，只要 \(\mathrm{Sp}(x)\) 非空。

\begin{enumerate}
\def\labelenumi{\arabic{enumi})}
\setcounter{enumi}{4}
\item
  设代数 A 非零。设 \(x \in A\) 是一个幂零元。设 n 是使 \(x^{n} = 0\) 的一个整数。\(x^{n}\) 的谱化为 0，从而由注 4，x 的谱亦然。
\item
  设 A 与 B 是 K 上的幺代数，\(\varphi\colon\mathrm{A}\to\mathrm{B}\) 是一个保单位元态射。对每个 \(x\in\mathrm{A}\)，有 \(\mathrm{Sp}_{\mathrm{B}}(\varphi(x))\subset\mathrm{Sp}_{\mathrm{A}}(x)\)。
\item
  设 A 是一个幺代数，R 是其根 (A, VIII, p.~150, 定义 2)，\(\varphi\) 是 A 到 \(B = A/R\) 上的典范态射。若 \(x \in A\)，则有 \(\mathrm{Sp}_{\mathrm{B}}(\varphi(x)) = \mathrm{Sp}_{\mathrm{A}}(x)\)。事实上，只需证明：若 \(\varphi(x)\) 在 B 中可逆，则 \(x\) 在 A 中可逆。而今，若 \(y \in A\) 使得 \(\varphi(x)\varphi(y) = \varphi(y)\varphi(x) = \varphi(e)\)，则有 \(xy \in e + R\)，\(yx \in e + R\)，因而 \(xy\) 与 \(yx\) 均可逆 (A, VIII, p.~151, 定理 1)，从而 \(x\) 可逆。特别地，若 \(x \in R\)，则我们有 \(\mathrm{Sp}_{\mathrm{A}}(x) = \{0\}\)，只要 \(A \neq \{0\}\)。
\item
  设 \((\mathrm{B}_{i})_{i\in\mathrm{I}}\) 是一族幺代数，其中 \(\mathrm{B}_{i}=(\mathrm{A}_{i},e_{i})\)（\(i\in I\)）。令 \(A=\prod_{i}A_{i}, e=(e_{i})_{i\in\mathrm{I}}\)。则 \((\mathrm{A},e)\) 是一个幺代数，称为诸 \(B_{i}\) 的积。若 \(x=(x_{i})_{i\in\mathrm{I}}\in\mathrm{A}\)，则有 \(\mathrm{Sp}_{\mathrm{A}}(x)=\bigcup_{i}\mathrm{Sp}_{\mathrm{A}_{i}}(x_{i})\)。
\end{enumerate}

例. --- 1) 设 X 是一个集合，A = K \(^{X}\) 是定义在 X 上而取值于 K 的函数所成的代数。A 的元素 f 的谱就是 f 的值集。

\begin{enumerate}
\def\labelenumi{\arabic{enumi})}
\setcounter{enumi}{1}
\tightlist
\item
  设 A 是 K 上有限秩的幺代数。要使 \(x \in A\) 可逆，必须且只需从 A 到 A 的线性映射 \(y \mapsto xy\) 的行列式非零。由此推出，x 的谱是 x 的特征多项式的根集 (A, III, p.~110)。因此，若 A 是 K 上有限维向量空间 V 的自同态代数，则 x 的谱就是 x 的特征值集。当 V 为无限维时，情况并非总是如此（参见 I, p.~153, 习题 2）。
\end{enumerate}

